\documentclass[
    language=english,
    doctype=video-script,
    institution=none,
]{../../omnilatex}

\RequirePackage{config/document-settings}

\begin{document}
\maketitle

\section*{VIDEO SCRIPT}

\textbf{Channel:} Tech Simplified\\
\textbf{Video:} ``Why Your Phone Battery Dies So Fast (And How to Fix It)''\\
\textbf{Host:} Jamie Chen\\
\textbf{Target Length:} 12 minutes\\
\textbf{Published:} June 12, 2024

\section*{VIDEO METADATA}

\begin{table}[h]
\centering
\begin{tabular}{lr}
\toprule
\textbf{Title} & Why Your Phone Battery Dies So Fast (And How to Fix It) \\
\textbf{Thumbnail} & Close-up of phone at 1\% battery, red icon, host's shocked face \\
\textbf{Tags} & phone battery, battery life, phone tips, tech tips, smartphone \\
\textbf{Category} & Science \& Technology \\
\textbf{CTA} & Like, subscribe, comment your battery life tips \\
\bottomrule
\end{tabular}
\caption{Video metadata}
\end{table}

\section*{SCRIPT}

\subsection*{HOOK (0:00--0:15)}

\textbf{B-ROLL:} Close-up of phone screen showing 1\% battery. Red warning icon. Phone dies.

\begin{quote}
JAMIE (V.O.): Your phone battery is dying faster than it should. And it's not because your battery is broken. It's because of five settings you probably don't even know exist. Let me show you how to fix them.
\end{quote}

\textbf{B-ROLL:} Quick montage of phone settings screens, battery graphs, host holding phone.

\begin{quote}
JAMIE (ON CAMERA): Hey everyone, welcome back to Tech Simplified. I'm Jamie Chen, and today we're talking about the number one complaint I hear from people about their phones: battery life. Stay until the end, because setting number five is the one that makes the biggest difference, and almost nobody knows about it.
\end{quote}

\textbf{CTA:} If you find this helpful, hit that like button. It really helps the channel.

\subsection*{SEGMENT 1: THE PROBLEM (0:15--2:00)}

\textbf{B-ROLL:} Infographic showing average screen-on time by phone model.

\begin{quote}
JAMIE (V.O.): Here's the thing: modern smartphones have bigger batteries than ever before. The iPhone 15 Pro Max has a 4,441 milliamp-hour battery. The Samsung Galaxy S24 Ultra has a 5,000 milliamp-hour battery. These are massive batteries. So why do they die so fast?
\end{quote}

\textbf{B-ROLL:} Animation showing battery drain graph over a typical day.

\begin{quote}
JAMIE (V.O.): The answer is software. Specifically, the apps and settings that are running in the background, consuming power even when you're not using your phone. The average smartphone has about 80 apps installed. Of those, maybe you use 30 regularly. The rest? They're sitting there, draining your battery, sending notifications, and tracking your location.
\end{quote}

\textbf{JAMIE (ON CAMERA):} So let's fix it. Here are the five settings that are killing your battery, and how to change them.

\subsection*{SEGMENT 2: SETTING 1 -- BACKGROUND APP REFRESH (2:00--3:30)}

\textbf{B-ROLL:} Screen recording showing how to disable background app refresh.

\begin{quote}
JAMIE (V.O.): Setting number one: background app refresh. This is the setting that allows apps to check for updates and new content even when you're not using them. It sounds useful, but in practice, it means that every app on your phone is constantly waking up, using your battery, and using your data.
\end{quote}

\textbf{JAMIE (ON CAMERA):} Here's how to turn it off. On iPhone, go to Settings, then General, then Background App Refresh. You can turn it off entirely, or you can choose which apps are allowed to refresh in the background. I recommend turning it off for everything except messaging apps and your calendar.

On Android, it's similar. Go to Settings, then Apps, then select an app. Look for ``Allow background activity'' and turn it off for apps that don't need it.

\textbf{B-ROLL:} Before/after battery usage graph showing improvement.

\begin{quote}
JAMIE (V.O.): I made this change on my phone last month, and my screen-on time increased by about 45 minutes per day. That's significant.
\end{quote}

\subsection*{SEGMENT 3: SETTING 2 -- LOCATION SERVICES (3:30--5:00)}

\textbf{B-ROLL:} Screen recording showing location services settings.

\begin{quote}
JAMIE (V.O.): Setting number two: location services. This is a big one. Most apps want to know your location all the time. Some of them need it---maps, ride-sharing, weather. But a lot of them don't. Why does your flashlight app need to know where you are? Why does your calculator need your location?
\end{quote}

\textbf{JAMIE (ON CAMERA):} Go to Settings, then Privacy, then Location Services. You'll see a list of every app that has requested your location. For each one, you can choose: Never, While Using the App, or Always. I recommend setting everything to ``While Using the App'' or ``Never'' except for apps that genuinely need continuous location access.

Also, turn off ``Precise Location'' for apps that don't need it. Most apps don't need to know exactly where you are. City-level is fine.

\textbf{B-ROLL:} Graphic showing battery savings from location services changes.

\begin{quote}
JAMIE (V.O.): Location services are one of the biggest battery drains, especially if you have apps set to ``Always.'' Turning this off can save you 10--15\% of your daily battery usage.
\end{quote}

\subsection*{SEGMENT 4: SETTING 3 -- DISPLAY SETTINGS (5:00--6:30)}

\textbf{B-ROLL:} Side-by-side comparison of display settings.

\begin{quote}
JAMIE (V.O.): Setting number three: your display. The screen is the single biggest battery drain on your phone. And there are two settings that make a huge difference: brightness and refresh rate.
\end{quote}

\textbf{JAMIE (ON CAMERA):} First, brightness. Turn on auto-brightness. I know some people prefer manual control, but auto-brightness will save you significant battery because it adjusts to your environment. If you're in a dark room, there's no reason to have your screen at 80\% brightness.

Second, refresh rate. If your phone has a 120Hz display---and a lot of modern phones do---switch it to 60Hz. Yes, 120Hz looks smoother. But it uses significantly more battery. Go to Settings, then Display, then Refresh Rate, and select 60Hz or ``Standard.''

\textbf{B-ROLL:} Visual comparison of 60Hz vs 120Hz scrolling.

\begin{quote}
JAMIE (V.O.): You'll barely notice the difference in daily use, but your battery will notice.
\end{quote}

\subsection*{SEGMENT 5: SETTING 4 -- NOTIFICATIONS (6:30--8:00)}

\textbf{B-ROLL:} Screen recording showing notification settings.

\begin{quote}
JAMIE (V.O.): Setting number four: notifications. Every time your phone gets a notification, the screen lights up, the speaker plays a sound, and the vibration motor fires. All of that uses battery. And most notifications? They're not important.
\end{quote}

\textbf{JAMIE (ON CAMERA):} Go to Settings, then Notifications. You'll see every app that's allowed to send you notifications. Go through this list and be honest with yourself: which of these apps actually need to notify you in real time? Your bank? Maybe. Your calendar? Yes. That game you played once six months ago? No.

Turn off notifications for everything except the apps that genuinely need to reach you. You'll be amazed at how much quieter your phone becomes---and how much longer your battery lasts.

\textbf{B-ROLL:} Graphic showing battery savings from notification management.

\begin{quote}
JAMIE (V.O.): The average phone receives 46 notifications per day. Each one uses a tiny amount of battery, but it adds up. Reducing your notifications by half can save you 5--10\% of your daily battery.
\end{quote}

\subsection*{SEGMENT 6: SETTING 5 -- THE BIG ONE (8:00--10:00)}

\textbf{B-ROLL:} Dramatic zoom on host's face.

\begin{quote}
JAMIE (ON CAMERA): Okay. Setting number five. This is the one that makes the biggest difference, and almost nobody knows about it. Ready?
\end{quote}

\begin{quote}
JAMIE (V.O.): It's called ``Optimized Battery Charging,'' and it's a setting that actually \textit{saves} your battery long-term by changing how your phone charges.
\end{quote}

\textbf{JAMIE (ON CAMERA):} Here's the problem: when you charge your phone to 100\% every night and leave it plugged in all night, you're actually degrading the battery faster. Lithium-ion batteries---which is what your phone uses---degrade when they're held at high charge levels for extended periods.

Optimized Battery Charging solves this by learning your charging patterns and holding the charge at 80\% until just before you normally unplug. So if you plug in at 11 PM and unplug at 7 AM, it will charge to 80\%, stop, and then finish charging to 100\% around 6:30 AM.

\textbf{B-ROLL:} Screen recording showing how to enable optimized battery charging.

\begin{quote}
JAMIE (V.O.): On iPhone, go to Settings, then Battery, then Battery Health and Charging, and turn on ``Optimized Battery Charging.'' On Android, go to Settings, then Battery, then Battery Care, and select ``Adaptive'' or ``Optimized.''
\end{quote}

\begin{quote}
JAMIE (ON CAMERA): This setting doesn't just improve your daily battery life---it extends the overall lifespan of your battery by months or even years. It's the single most impactful change you can make.
\end{quote}

\subsection*{SEGMENT 7: RECAP AND CTA (10:00--12:00)}

\textbf{B-ROLL:} Quick recap graphics for each setting.

\begin{quote}
JAMIE (ON CAMERA): So let's recap. The five settings that are killing your battery:
\end{quote}

\begin{enumerate}
    \item Background App Refresh --- turn it off for most apps
    \item Location Services --- set to ``While Using'' or ``Never''
    \item Display --- enable auto-brightness, reduce refresh rate
    \item Notifications --- turn off everything non-essential
    \item Optimized Battery Charging --- turn it on
\end{enumerate}

\begin{quote}
JAMIE (ON CAMERA): If you make all five of these changes, you should see a significant improvement in your battery life. I'm talking 20--30\% more screen-on time, which for most people means getting through the entire day without needing to charge.

Now I want to hear from you. What are your best battery-saving tips? Drop them in the comments below. I read every single one, and I might feature my favorites in a future video.

If you found this video helpful, please give it a thumbs up. It really helps the channel grow. And if you haven't already, hit that subscribe button and the bell icon so you don't miss future videos.

Thanks for watching. I'll see you in the next one.
\end{quote}

\textbf{B-ROLL:} End screen with subscribe button and suggested videos.

\begin{quote}
\textbf{[END]}
\end{quote}

\section*{DESCRIPTION BOX}

\begin{quote}
Why does your phone battery die so fast? In this video, I show you the 5 settings that are draining your battery---and how to fix them in under 10 minutes.

These tips work for iPhone and Android, and they can increase your battery life by 20--30\%.

TIMESTAMPS:\\
0:00 - Why your battery dies fast\\
2:00 - Setting 1: Background App Refresh\\
3:30 - Setting 2: Location Services\\
5:00 - Setting 3: Display Settings\\
6:30 - Setting 4: Notifications\\
8:00 - Setting 5: Optimized Battery Charging\\
10:00 - Recap

SUBSCRIBE: [link]\\
MERCH: [link]\\
INSTAGRAM: @techsimplified\\
TWITTER: @techsimplified

#techtips #batterylife #smartphone
\end{quote}

\end{document}
